\chapter{ต้นไม้ตัดสินใจและการสุ่มป่าไม้}
ต้นไม้ตัดสินใจ (Decision Tree) เป็นอัลกอริทึมการเรียนรู้ของเครื่องที่มีความสำคัญและได้รับความนิยมอย่างกว้างขวางในการแก้ปัญหาทั้งการจำแนกประเภท (Classification) และการทำนายค่า (Regression) เนื่องจากมีโครงสร้างที่เข้าใจง่าย สามารถตีความผลได้ชัดเจน และไม่ต้องการการปรับแต่งข้อมูลมากนัก อย่างไรก็ตาม ต้นไม้ตัดสินใจเดี่ยวมักมีปัญหา Overfitting และมีความไวต่อการเปลี่ยนแปลงของข้อมูล ซึ่งนำไปสู่การพัฒนาเทคนิค Ensemble Learning โดยเฉพาะการสุ่มป่าไม้ (Random Forest) ที่เป็นการรวมต้นไม้ตัดสินใจหลายต้นเข้าด้วยกันผ่านกระบวนการ Bagging และการสุ่มเลือกตัวแปรย่อย (Feature Sampling) เพื่อเพิ่มความแม่นยำและลดปัญหา Overfitting ทำให้ Random Forest กลายเป็นหนึ่งในอัลกอริทึมที่ทรงพลังและใช้งานได้จริงในปัญหาหลากหลายสาขา ตั้งแต่การวิเคราะห์ข้อมูลทางการแพทย์ การเงิน ไปจนถึงการจำแนกภาพและการประมวลผลภาษาธรรมชาติ ในบทนี้จะกล่าวถึงโครงสร้างและหลักการทำงานของต้นไม้ตัดสินใจ กระบวนการแบ่งข้อมูลโดยใช้แนวคิด Entropy และ Information Gain Ratio การตัดแต่งกิ่ง (Pruning) เพื่อป้องกัน Overfitting รวมถึงหลักการของ Random Forest และเทคนิค Ensemble Learning ที่นำไปสู่ประสิทธิภาพที่ดีขึ้น


\section{ต้นไม้ตัดสินใจ (Decision Trees)}
\subsection{โครงสร้างและหลักการทำงาน}
ต้นไม้ตัดสินใจ (Decision Tree) เป็นโครงสร้างข้อมูลที่ประกอบด้วยโหนด (Node) และกิ่ง (Branch) โดยโหนดภายใน (Internal Node) แทนการทดสอบเงื่อนไขบนตัวแปร แต่ละกิ่งแทนผลลัพธ์ของการทดสอบ และโหนดลีฟ (Leaf Node) แทนค่าคำตอบ (Class Label หรือ Regression Value)

\begin{exmp}
\label{exmp:decision_tree_structure}
พิจารณาตัวอย่างการตัดสินใจอนุมัติสินเชื่อ ต้นไม้ตัดสินใจมีโครงสร้างดังนี้

\begin{figure}[h!]
\centering
%\includegraphics[width=0.8\textwidth]{decision_tree_loan_example.png}
\caption{ตัวอย่างต้นไม้ตัดสินใจสำหรับการอนุมัติสินเชื่อ}
\label{fig:loan_decision_tree}
\end{figure}

โหนดราก (Root Node): ทดสอบ \texttt{รายได้ $\geq$ 50,000 บาท/เดือน?}
\begin{itemize}
\item ถ้า \textbf{ใช่} $\rightarrow$ ไปทดสอบ \texttt{มีหลักประกัน?}
  \begin{itemize}
  \item ถ้ามีหลักประกัน $\rightarrow$ อนุมัติ
  \item ถ้าไม่มี $\rightarrow$ ตรวจสอบประวัติเครดิต
  \end{itemize}
\item ถ้า \textbf{ไม่} $\rightarrow$ ไปทดสอบ \texttt{ประวัติเครดิตดี?}
  \begin{itemize}
  \item ถ้าดี $\rightarrow$ อนุมัติแบบมีเงื่อนไข
  \item ถ้าไม่ดี $\rightarrow$ ไม่อนุมัติ
  \end{itemize}
\end{itemize}

จากต้นไม้นี้ ผู้กู้ที่มีรายได้สูงและมีหลักประกันจะได้รับการอนุมัติ ในขณะที่ผู้กู้ที่มีรายได้ต่ำและประวัติเครดิตไม่ดีจะถูกปฏิเสธ
\end{exmp}

\subsection{เอนโทรปี}
เอนโทรปี (Entropy) เป็นมาตรวัดความไม่บริสุทธิ์ (Impurity) ของชุดข้อมูล ใช้ในการวัดว่าชุดข้อมูลมีความหลากหลายมากน้อยเพียงใด สูตรของเอนโทรปีคือ

\[
\text{Entropy}(S) = -\sum_{i=1}^{c} p_i \log_2(p_i)
\]

โดยที่ $p_i$ คือสัดส่วนของตัวอย่างในคลาส $i$ และ $c$ คือจำนวนคลาสทั้งหมด

\begin{exmp}
\label{exmp:entropy_calculation}
พิจารณาชุดข้อมูลลูกค้า 10 ราย ที่ต้องการสินเชื่อ โดยมีผลลัพธ์ดังนี้
\begin{itemize}
\item อนุมัติ: 7 ราย ($p_1 = 0.7$)
\item ไม่อนุมัติ: 3 ราย ($p_2 = 0.3$)
\end{itemize}

ค่าเอนโทรปีของชุดข้อมูลนี้คำนวณได้ดังนี้

\[
\text{Entropy}(S) = -(0.7)\log_2(0.7) - (0.3)\log_2(0.3) \approx 0.610 + 0.521 = 1.131 \text{ bits}
\]

หากชุดข้อมูลมีเพียงคลาสเดียว (ทั้งหมดอนุมัติหรือไม่อนุมัติ) ค่าเอนโทรปีจะเท่ากับ 0 ซึ่งหมายถึงความบริสุทธิ์สมบูรณ์ ในทางกลับกัน หากชุดข้อมูลมีการกระจาย 50-50 ค่าเอนโทรปีจะเท่ากับ 1 bit ซึ่งเป็นค่าสูงสุด
\end{exmp}

\subsection{เกรนความรู้และอัตราเกรนความรู้}
เกรนความรู้ (Information Gain) วัดการลดลงของเอนโทรปีหลังจากแบ่งข้อมูลตามตัวแปรใดตัวแปรหนึ่ง สูตรคือ

\[
\text{IG}(S, A) = \text{Entropy}(S) - \sum_{v \in \text{Values}(A)} \frac{|S_v|}{|S|} \text{Entropy}(S_v)
\]

อัตราเกรนความรู้ (Gain Ratio) เป็นการปรับปรุง Information Gain เพื่อลดอคติที่ชอบตัวแปรที่มีค่าหลายค่า

\[
\text{GR}(S, A) = \frac{\text{IG}(S, A)}{\text{SplitInformation}(S, A)}
\]

\begin{exmp}
\label{exmp:information_gain}
จากข้อมูลตัวอย่าง 10 รายในตัวอย่าง \ref{exmp:entropy_calculation} พิจารณาการแบ่งข้อมูลตามตัวแปร "ประเภทรายได้" ซึ่งมี 2 ค่า คือ "ประจำ" และ "Freelance"

\begin{table}[h!]
\centering
\caption{ข้อมูลการแบ่งตามประเภทรายได้}
\label{tab:income_split}
\begin{tabular}{|c|c|c|c|}
\hline
\textbf{ประเภทรายได้} & \textbf{จำนวน} & \textbf{อนุมัติ} & \textbf{ไม่อนุมัติ} \\ \hline
ประจำ & 6 & 5 & 1 \\ \hline
Freelance & 4 & 2 & 2 \\ \hline
\end{tabular}
\end{table}

ค่าเอนโทรปีหลังแบ่ง:
\begin{itemize}
\item กลุ่มประจำ: $\text{Entropy} = -(5/6)\log_2(5/6) - (1/6)\log_2(1/6) \approx 0.650$ bits
\item กลุ่ม Freelance: $\text{Entropy} = -(2/4)\log_2(2/4) - (2/4)\log_2(2/4) = 1.0$ bit
\end{itemize}

ค่า Information Gain:
\[
\text{IG} = 1.131 - \left(\frac{6}{10}(0.650) + \frac{4}{10}(1.0)\right) \approx 1.131 - 0.790 = 0.341 \text{ bits}
\]

ค่า Information Gain นี้แสดงว่าการแบ่งตามประเภทรายได้ช่วยลดความไม่แน่นอนลงได้ 0.341 bits
\end{exmp}

\section{การสุ่มป่าไม้ (Random Forest)}
\subsection{แนวคิด Ensemble Learning}
Ensemble Learning เป็นเทคนิคที่รวมการทำนายจากตัวแบบหลายตัวเพื่อให้ได้ผลลัพธ์ที่ดีกว่าตัวแบบใดตัวแบบหนึ่ง โดยอาศัยหลักการที่ว่าความผิดพลาดของตัวแบบแต่ละตัวมักจะไม่สัมพันธ์กัน (Uncorrelated Errors)

\begin{exmp}
\label{exmp:ensemble_voting}
พิจารณาการจำแนกประเภทอีเมล 3 แบบ (Spam, Not Spam, Unknown) โดยใช้ตัวแบบ 5 ตัวใน Ensemble ผลลัพธ์การโหวตดังนี้

\begin{itemize}
\item ตัวแบบ 1: Spam
\item ตัวแบบ 2: Not Spam
\item ตัวแบบ 3: Spam
\item ตัวแบบ 4: Spam
\item ตัวแบบ 5: Unknown
\end{itemize}

ผลลัพธ์สุดท้าย (Majority Voting): Spam (3 จาก 5 โหวต)

หากใช้ต้นไม้ตัดสินใจเพียงต้นเดียว อาจให้ผลลัพธ์ที่ไม่ถูกต้อง แต่การรวมการโหวตจากหลายต้นช่วยลดความเสี่ยงของความผิดพลาดจากต้นใดต้นหนึ่ง
\end{exmp}

\subsection{Bagging และ Bootstrap Sampling}
Bagging (Bootstrap Aggregating) เป็นวิธีการสร้างตัวแบบ Ensemble โดยการสุ่มข้อมูลย่อย (Bootstrap Samples) และฝึกตัวแบบแต่ละตัวบนข้อมูลย่อยเหล่านั้น จากนั้นรวมผลลัพธ์ด้วยการโหวตหรือการเฉลี่ย

\begin{exmp}
\label{exmp:bagging_example}
จากชุดข้อมูล Training ขนาด $n = 10$ การสร้าง Bootstrap Sample ที่ 1 อาจได้ $\{x_3, x_1, x_7, x_1, x_5, x_9, x_1, x_2, x_4, x_7\}$ (สังเกตว่าบางรายการถูกเลือกซ้ำ บางรายการไม่ถูกเลือก)

Bootstrap Sample ที่ 2 อาจได้ $\{x_2, x_5, x_8, x_1, x_3, x_6, x_10, x_4, x_7, x_9\}$

การสุ่มแบบนี้ช่วยให้แต่ละต้นไม้เห็นข้อมูลที่แตกต่างกันเล็กน้อย ทำให้ความผิดพลาดของต้นไม้แต่ละต้นไม่สัมพันธ์กันอย่างสมบูรณ์
\end{exmp}

\subsection{ความสำคัญของตัวแปร (Feature Importance)}
ความสำคัญของตัวแปร (Feature Importance) เป็นการวัดว่าแต่ละตัวแปรมีส่วนร่วมในการลดความไม่บริสุทธิ์ (Impurity) ของโหนดในต้นไม้มากน้อยเพียงใด โดยทั่วไปคำนวณจากผลรวมของ Information Gain ที่ตัวแปรนั้นถูกใช้ในการแบ่งโหนด รวมถึงน้ำหนักตามจำนวนตัวอย่างในโหนดนั้น

\begin{exmp}
\label{exmp:feature_importance}
จากการฝึก Random Forest บนข้อมูลการอนุมัติสินเชื่อ ค่า Feature Importance ของตัวแปรต่าง ๆ มีดังนี้

\begin{table}[h!]
\centering
\caption{ค่า Feature Importance จาก Random Forest}
\label{tab:feature_importance}
\begin{tabular}{|c|c|}
\hline
\textbf{ตัวแปร} & \textbf{Feature Importance} \\ \hline
รายได้ & 0.42 \\ \hline
ประวัติเครดิต & 0.31 \\ \hline
อายุงาน & 0.15 \\ \hline
หลักประกัน & 0.09 \\ \hline
อื่น ๆ & 0.03 \\ \hline
\end{tabular}
\end{table}

จากตารางจะเห็นได้ว่า "รายได้" และ "ประวัติเครดิต" เป็นตัวแปรที่มีความสำคัญมากที่สุดในการตัดสินใจอนุมัติสินเชื่อ ตัวแปรเหล่านี้ควรได้รับความสนใจเป็นพิเศษในการวิเคราะห์และการเก็บข้อมูล
\end{exmp}



\newpage
%\RefChapter